% Self-contained contribution; standard amsmath/amssymb/amsthm environments.
% Suggested bibliography keys: incoming-shifted, incoming-convolution,
% incoming-closure, dlmf-hurwitz.  All statements use ambient x cutoffs.
\section{Dilation and unequal-frequency correlations}
\label{mld:dil:section}\label{mld:sec:dilation}

Two explicit questions in the incoming reports concern the effect of
integer dilation on canonical periodic finite parts, and products with
unequal integer dilations.  They occur in the further-research sections
of the shifted-Hurwitz and convolution-calculus deliveries.  This section
answers those questions for every nonnegative Stieltjes index and every
nonnegative parameter derivative order, provided the two sets of
singular points are disjoint.  The location of the surviving ordinary
special-function values is determined by the reduced dilation ratio;
their coefficients also remember the physical coordinate used for
the finite part.

\subsection{Conventions and the elementary contact identity}
Write $\mathbb T=\mathbb R/\mathbb Z$, and use normalized Lebesgue
measure.  The spectral convention is
\[
 \zeta(1-z,x)=-\frac1z+
       \sum_{n\ge0}\gamma_n(x)\frac{z^n}{n!},
 \qquad \gamma_0(x)=-\psi(x),
 \qquad 0<x<1.
\]
For a smooth periodic test function $\phi$, let $T_{n,r}$ be the
constant-term extension of $\gamma_n^{(r)}(x)$ from $(0,1)$:
\[
 \langle T_{n,r},\phi\rangle
 =\operatorname{CT}_{\epsilon\downarrow0}
      \int_\epsilon^1\gamma_n^{(r)}(x)\phi(x)\,dx.
\]
Here $\operatorname{CT}$ retains the constant term after all negative
powers of $\epsilon$ multiplied by powers of $\log\epsilon$, and all
positive powers of $\log\epsilon$, have been removed.  No finite
point-supported term is added.  Thus, with $G_n=T_{n,0}$,
\[
 \langle G_n,\phi\rangle
 =\int_0^1\left\{\gamma_n(x)\phi(x)
                  -\frac{\log^n x}{x}\phi(0)\right\}\,dx.
\]
The shift identity $\gamma_n(x)=x^{-1}\log^n x+\gamma_n(1+x)$
proves convergence and the equality of these definitions.

Put
\begin{equation}
 p_r(z)=\prod_{j=1}^r\left(1-\frac zj\right),\qquad
 h_{r,n}=(-1)^n n!e_{n+1}(1,1/2,\ldots,1/r),
 \label{mld:dil:p-h}
\end{equation}
where the empty product is one and elementary symmetric functions of
degree greater than $r$ are zero.  The undilated contact identity, proved
in the incoming convolution and shifted-Hurwitz reports
\cite{IncomingConvolution,IncomingShifted}, is
\begin{equation}
 D^rG_n=T_{n,r}-h_{r,n}\delta_0^{(r)}.
 \label{mld:dil:unit-contact}
\end{equation}
For completeness, its origin can be read directly from the meromorphic
periodic Hurwitz family.  If
\[
 \widehat R_z(0)=0,\qquad
 \widehat R_z(k)=(2\pi i k)^{-z}\quad(k\ne0),
 \qquad Z(1-z)=\Gamma(z)R_z,
\]
then
\[
 Z(1-z)=\frac{\delta_0-1}{z}
            +\sum_{n\ge0}G_n\frac{z^n}{n!}.
\]
For $r\ge1$, the local singular part of the pointwise $r$th derivative
is $(-1)^r r!p_r(z)x^{z-r-1}$.  The meromorphic extension of
$x^{z-r-1}$ has residue $(-1)^r\delta_0^{(r)}/r!$ at $z=0$.
Consequently its polar contribution is
$p_r(z)\delta_0^{(r)}/z$.  Comparing coefficients with
$D^rZ(1-z)$ gives \eqref{mld:dil:unit-contact}, since
$n![z^{n+1}]p_r(z)=-h_{r,n}$.
The case $r=0$ is immediate.  This derivation also specifies all signs.
The coefficients can equally be written through generalized harmonic
numbers: for $|z|<1$,
\[
 p_r(z)=\exp\!\left(-\sum_{j\ge1}\frac{H_r^{(j)}}jz^j\right),
 \qquad H_r^{(j)}=\sum_{\ell=1}^r\ell^{-j}.
\]
For example $e_1=H_r$ and $e_2=(H_r^2-H_r^{(2)})/2$.

\subsection{The exact dilation anomaly}
For a positive integer $q$, let $U_q$ be pullback by the covering map
$x\mapsto qx$:
\[
 \langle U_qT,\phi\rangle
 =\left\langle T,\frac1q\sum_{j=0}^{q-1}
                         \phi\!\left(\frac{x+j}{q}\right)\right\rangle.
\]
For an ordinary function this is $U_qf(x)=f(qx)$.  Write
\[
 \Delta_q=\sum_{j=0}^{q-1}\delta_{j/q},\qquad
 U_q\delta_0^{(r)}=q^{-r-1}\Delta_q^{(r)}.
\]
Let $X_{q,n,r}$ be the constant-term extension of
$\gamma_n^{(r)}(\{qx\})$ using the \emph{ambient coordinate $x$}:
at each grid point $j/q$, delete the right interval
$(j/q,j/q+\epsilon)$ and take the constant term as $\epsilon$ tends to
zero. The transported finite part $U_qT_{n,r}$ takes the constant term
in a source cutoff $\eta$, corresponding to ambient cutoff $\eta/q$.
The ambient finite part instead takes the constant term in $\epsilon$
after the source lower limit becomes $q\epsilon$. The variable in which
the constant term is extracted is part of the normalization.

\begin{theorem}[All-order finite-part dilation law]
\label{mld:dil:scale-thm}
For $q\ge1$ and $n,r\ge0$, define the explicit polynomial
\begin{equation}
 c_{n,r}(L)=n![z^{n+1}]p_r(z)(e^{Lz}-1)
 =n!\sum_{j=0}^{\min(r,n)}
       \frac{(-1)^j e_j(1,1/2,\ldots,1/r)}{(n+1-j)!}
       L^{n+1-j}.
 \label{mld:dil:c}
\end{equation}
Then
\begin{equation}
 \boxed{X_{q,n,r}=U_qT_{n,r}
       -q^{-r-1}c_{n,r}(\log q)\Delta_q^{(r)}.}
 \label{mld:dil:scale-law}
\end{equation}
In particular $c_{n,0}(L)=L^{n+1}/(n+1)$ and $c_{0,r}(L)=L$.
\end{theorem}

\begin{proof}
Near a grid point use $t=x-j/q>0$.  The singular part of the generating
function for the argument derivatives is
\[
 (-1)^r r!p_r(z)(qt)^{z-r-1}
 =(-1)^r r!q^{-r-1}p_r(z)e^{z\log q}t^{z-r-1}.
\]
After pairing with a test function, only its Taylor term of degree $r$
produces a pole at $z=0$.  Its contribution is
\[
 q^{-r-1}\frac{p_r(z)e^{z\log q}}z\delta_{j/q}^{(r)}.
\]
The pullback of the source-coordinate family instead has polar
contribution $q^{-r-1}p_r(z)\delta_{j/q}^{(r)}/z$.
The regular Taylor coefficients left after subtracting these complete
polar contributions are respectively the ambient-coordinate constant
terms and the pulled-back constant terms.  Their difference is therefore
the negative of
\[
 q^{-r-1}p_r(z)\frac{e^{z\log q}-1}{z}\delta_{j/q}^{(r)}.
\]
Extract $z^n/n!$ and sum over the grid.  When $r=0$, the scalar spectral
pole $-1/z$ occurs in both calculations and cancels.  Subtracting the
test-function Taylor polynomial through degree $r$ makes every
remaining local integral holomorphic near $z=0$, which justifies the
coefficient comparison.  Expanding the finite product gives
\eqref{mld:dil:c}.
\end{proof}

\begin{corollary}[Differentiation and dilation with fixed ambient scale]
Define
\begin{equation}
 b_{r,n}(L)=\sum_{j=0}^n\binom nj L^{n-j}h_{r,j}.
 \label{mld:dil:b}
\end{equation}
Then
\begin{equation}
 \boxed{D^rX_{q,n,0}=q^rX_{q,n,r}
          -\frac{b_{r,n}(\log q)}q\Delta_q^{(r)}.}
 \label{mld:dil:commutator}
\end{equation}
For every $n\ge0$ this specializes to
\begin{equation}
 DX_{q,n,0}=qX_{q,n,1}
             -\frac{(\log q)^n}{q}\Delta_q',
 \label{mld:dil:first-commutator}
\end{equation}
where $(\log1)^0=1$.
\end{corollary}
\begin{proof}
Differentiate \eqref{mld:dil:scale-law} for $r=0$, use
$D^rU_q=q^rU_qD^r$ and \eqref{mld:dil:unit-contact}, and compare with
\eqref{mld:dil:scale-law} at general $r$.  The remaining scalar identity is
\[
 b_{r,n}(L)=h_{r,n}+\frac{L^{n+1}}{n+1}-c_{n,r}(L)
          =n![z^{n+1}](1-p_r(z))e^{Lz}.
\]
For $r=1$, only $h_{1,0}=1$ is nonzero.
\end{proof}

The undilated anomaly vanishes when $n\ge r$.  Equation
\eqref{mld:dil:first-commutator} shows that this vanishing generally fails
after ambient-coordinate dilation: already $n=r=1$ gives
$DX_{q,1,0}=qX_{q,1,1}-(\log q)\Delta_q'/q$.
This is an extension of the incoming contact theorem, not a correction
to its correctly normalized undilated statement.

The scale corrections are compatible with successive dilations.  The
following polynomial identities hold for all $L,M\in\mathbb C$:
\begin{align}
 c_{n,r}(L+M)&=c_{n,r}(M)
       +\sum_{j=0}^n\binom nj M^{n-j}c_{j,r}(L),\notag\\
 b_{r,n}(L+M)&=\sum_{j=0}^n\binom nj M^{n-j}b_{r,j}(L).
 \label{mld:dil:scale-composition}
\end{align}
Indeed, use $e^{(L+M)z}-1=(e^{Mz}-1)+e^{Mz}(e^{Lz}-1)$ in
\eqref{mld:dil:c}, and use the last generating expression in the preceding
proof for $b$.  Thus a two-step change of local scale gives the same
finite correction as the combined scale; no independent finite
constant is introduced by composing the operations.

\subsection{A multiplication theorem including the point masses}
Let $\tau_aT(x)=T(x+a)$ and $L=\log q$.
\begin{theorem}[Distributional multiplication law]
For $n\ge0$,
\begin{equation}
 \sum_{j=0}^{q-1}\tau_{j/q}G_n
 =q\sum_{i=0}^n\binom ni(-L)^{n-i}U_qG_i
     +\frac{(-L)^{n+1}}{n+1}(\Delta_q-q).
 \label{mld:dil:multiplication}
\end{equation}
More generally, for $r\ge1$,
\begin{equation}
 \sum_{j=0}^{q-1}\tau_{j/q}T_{n,r}
 =q^{r+1}\sum_{i=0}^n\binom ni(-L)^{n-i}U_qT_{i,r}
                  +c_{n,r}(-L)\Delta_q^{(r)}.
 \label{mld:dil:derivative-multiplication}
\end{equation}
\end{theorem}
\begin{proof}
The classical Hurwitz multiplication formula gives, as meromorphic
periodic distributions,
\[
 \sum_{j=0}^{q-1}\tau_{j/q}Z(1-z)
                  =q^{1-z}U_qZ(1-z).
\]
It can also be proved directly from the Fourier coefficients: the
sum annihilates frequencies not divisible by $q$.  Expanding the
distributional Laurent series proves \eqref{mld:dil:multiplication}; the
residue transforms as $\sum\tau_{j/q}(\delta_0-1)=\Delta_q-q$.

Differentiate the pointwise multiplication formula $r$ times and then
take ambient-coordinate constant terms.  Substitution of
\eqref{mld:dil:scale-law} reduces the additional coefficient to
\[
 -\sum_{i=0}^n\binom ni(-L)^{n-i}c_{i,r}(L)
 =n![z^{n+1}]p_r(z)(e^{-Lz}-1)=c_{n,r}(-L).
\]
This gives \eqref{mld:dil:derivative-multiplication}.
\end{proof}

\subsection{The undilated kernel and an explicit finite coefficient rule}
\label{mld:dil:undilated-algorithm}

The unequal-dilation theorem below uses a known undilated closure from
the incoming reports \cite{IncomingClosure,IncomingShifted}. We record its
generating function and derivation so that the resulting algorithm is
self-contained. These undilated statements are prior results within the
supplied material.

For $0<t<1$ and $\Re s,\Re v<1$, set
\[
 C(s,v;t)=\int_0^1\zeta(s,x)\zeta(v,\{x+t\})\dd x.
\]
The integral is ordinary: the two singularities are separated. Its two
equivalent continuations are
\begin{align}
 C(s,v;t)={}&\Gamma(1-s)\Gamma(1-v)(2\pi)^{s+v-2}
 \bigl\{e^{-\pi\ii(s-v)/2}\Li_{2-s-v}(e^{2\pi\ii t})\notag\\
 &\hspace{38mm}+e^{\pi\ii(s-v)/2}\Li_{2-s-v}(e^{-2\pi\ii t})\bigr\},
 \label{mld:dil:C-polylog}\\
 C(s,v;t)={}&\Gamma(s+v-1)
 \left\{\frac{\Gamma(1-s)}{\Gamma(v)}\zeta(s+v-1,t)
       +\frac{\Gamma(1-v)}{\Gamma(s)}\zeta(s+v-1,1-t)\right\}.
 \label{mld:dil:C-hurwitz}
\end{align}
All apparent singularities are interpreted in the complete expression.
To prove the first formula, start with $\Re s,\Re v<0$ and the absolutely
convergent Hurwitz Fourier series, whose nonzero coefficients are
\[
 \widehat\zeta_s(k)=\Gamma(1-s)(2\pi\ii k)^{s-1},\qquad k\ne0.
\]
The zero coefficient is zero. Integrating the product retains the
opposite frequencies and yields \eqref{mld:dil:C-polylog}. For the second
formula, insert the same Fourier expansion with spectral order $s+v-1$
on its right. Euler's gamma reflection identity reduces the coefficient
of the positive-argument polylogarithm by
\[
 \frac{\sin(\pi s)e^{-\pi\ii(s+v)/2}
       +\sin(\pi v)e^{\pi\ii(s+v)/2}}{\sin\pi(s+v)}
       =e^{-\pi\ii(s-v)/2}.
\]
The other coefficient is obtained by reversing the signs in the phases.
This calculation holds off the integer hyperplanes; analytic continuation
extends it across their removable singularities. The local decomposition
$\zeta(s,x)=x^{-s}+\zeta(s,1+x)$ justifies joint holomorphy of the
integral on $\Re s,\Re v<1$, and hence the continuation to that region.
This is the classical Fourier and product-integral mechanism
\cite{DLMFhurwitz,EspinosaMoll} in the circular-translation convention.

Define, near $(u,v)=(0,0)$,
\begin{align*}
 \mathcal A(u,v)&=
  \frac{\Gamma(1+u+v)\Gamma(1-v)}{\Gamma(1+u)},
 &\mathcal B(u,v)&=\mathcal A(v,u),\\
 \mathcal K(u,v;t)&=uv\,C(1+u,1+v;t).
\end{align*}
Equation~\eqref{mld:dil:C-hurwitz} gives
\begin{equation}
 \mathcal K(u,v;t)=
 -u\mathcal A(u,v)\zeta(1+u+v,1-t)
 -v\mathcal B(u,v)\zeta(1+u+v,t).
 \label{mld:dil:K-hurwitz}
\end{equation}
This function is holomorphic at the origin. Indeed,
$\mathcal A(u,-u)=\mathcal B(u,-u)=1$, so
$(u\mathcal A+v\mathcal B)/(u+v)$ has a removable denominator. Removing
the common Hurwitz pole in \eqref{mld:dil:K-hurwitz} leaves the negative of that
quotient plus analytic Stieltjes series.

For the canonical undilated finite part one consequently has the
terminating coefficient rule
\begin{equation}
 \boxed{I_{mn}(t)=(-1)^{m+n}m!n!
       [u^{m+1}v^{n+1}]\mathcal K(u,v;t).}
 \label{mld:dil:I-coefficient}
\end{equation}
Here $I_{mn}$ denotes the constant term of
$\int\gamma_m(x)\gamma_n(\{x+t\})\dd x$ with the two ambient cutoffs.
For clarity, the matching to a finite part can be justified before
extracting any coefficients. Put
$\mathcal H(u,x)=\zeta(1+u,x)-1/u$ and $b=1-t$.
The following locally convergent expression is analytic for small $u,v$:
\begin{align*}
 &\int_0^b\left\{
 x^{-1-u}[\mathcal H(v,t+x)-\mathcal H(v,t)]
       +\mathcal H(u,1+x)\mathcal H(v,t+x)\right\}\dd x\\
 &+\int_0^t\left\{
 y^{-1-v}[\mathcal H(u,b+y)-\mathcal H(u,b)]
       +\mathcal H(v,1+y)\mathcal H(u,b+y)\right\}\dd y\\
 &+\mathcal H(v,t)\frac{1-b^{-u}}u
       +\mathcal H(u,b)\frac{1-t^{-v}}v.
\end{align*}
The bracketed differences vanish linearly; both quotients have removable
values at zero. The coefficient of $u^mv^n$ is exactly
$(-1)^{m+n}I_{mn}(t)/(m!n!)$, by direct integration of the subtracted
logarithmic singularities. In a left spectral half-plane the displayed
expression is also
\[
 C(1+u,1+v;t)+\frac1{uv}
       +\frac{\mathcal H(v,t)}u+\frac{\mathcal H(u,b)}v.
\]
Using the axis values of \eqref{mld:dil:K-hurwitz}, it equals
\[
 \frac{\mathcal K(u,v;t)-\mathcal K(u,0;t)
            -\mathcal K(0,v;t)+\mathcal K(0,0;t)}{uv}.
\]
This proves \eqref{mld:dil:I-coefficient} with the specified normalization.

The coefficient calculation uses rational polynomial arithmetic and
positive integer zeta values only, because
\begin{equation}
 \log\mathcal A(u,v)=\sum_{j\ge2}\frac{\zeta(j)}j
       \left\{(-1)^j\bigl((u+v)^j-u^j\bigr)+v^j\right\}.
 \label{mld:dil:logA}
\end{equation}
For an explicit finite algorithm, let $\ell_j$ be the homogeneous summand
of degree $j$ here. The homogeneous terms of $\mathcal A$ satisfy
\[
 \mathcal A_0=1,\qquad\mathcal A_1=0,\qquad
 \mathcal A_d=\frac1d\sum_{j=2}^d j\ell_j\mathcal A_{d-j}.
\]
Only $d\le m+n+2$ can contribute to \eqref{mld:dil:I-coefficient}. This proves
finite linear closure in the single values $\gamma_j(t)$ and
$\gamma_j(1-t)$, with $j\le m+n+1$. Their top coefficients are
$1/(m+1)$ and $1/(n+1)$; the coefficients at order $m+n$ vanish since
$\mathcal A_1=0$. These are the incoming closure theorem and its
missing-order rule, now available as an explicit input to the new
dilation calculation.


\subsection{Unequal integer dilations: the full Stieltjes theorem}
Fix positive integers $p,q$, put
\begin{equation}
 d=\gcd(p,q),\qquad P=p/d,\qquad Q=q/d,\qquad
 \theta=\{Pa\},\qquad 0<\theta<1.
 \label{mld:dil:parameters}
\end{equation}
The last condition is exactly the condition that the singular grids of
$\gamma_m(\{px\})$ and $\gamma_n(\{qx+a\})$ do not intersect.
Indeed an intersection would give $qx+a\in\mathbb Z$ and
$px\in\mathbb Z$, equivalently $Pa\in\mathbb Z$ since $(P,Q)=1$.
Let
\[
 F_{mn}^{p,q}(\theta)
 =\operatorname{FP}_x\int_0^1
       \gamma_m(\{px\})\gamma_n(\{qx+a\})\,dx.
\]
The notation $\operatorname{FP}_x$ means the common ambient-coordinate
constant term specified above. One may equivalently take a separate
constant term in an independent ambient cutoff variable at each grid
point. Replacing those variables by $c_j\epsilon$ and taking a single
constant term would generally change the answer through local
$\log c_j$ terms.

We use the already established undilated correlation
\[
 I_{mn}(t)=\operatorname{FP}\int_0^1
                  \gamma_m(x)\gamma_n(\{x+t\})\,dx,
 \qquad 0<t<1.
\]
Its all-order single-Stieltjes closure and finite coefficient rule were
recalled in Section~\ref{mld:dil:undilated-algorithm}.  The formulas needed below are
\begin{align}
 I_{00}(t)&=\gamma_1(t)+\gamma_1(1-t)-2\zeta(2),
 \label{mld:dil:I00}\\
 I_{10}(t)&=\tfrac12\gamma_2(t)+\gamma_2(1-t)
        +\zeta(2)\{\gamma_0(t)+\gamma_0(1-t)\}-\zeta(3).
 \label{mld:dil:I10}
\end{align}
For $M\ge1$ define a normalized trace polynomial
\begin{equation}
 A_n(M,t)=\sum_{j=0}^n\binom nj(-\log M)^{n-j}\gamma_j(t)
                 -\frac{(-\log M)^{n+1}}{n+1}.
 \label{mld:dil:trace-polynomial}
\end{equation}
In particular $A_0(M,t)=\gamma_0(t)+\log M$.
The pointwise Hurwitz multiplication identity \cite[\S25.11]{DLMFhurwitz} proves
\[
 \frac1M\sum_{j=0}^{M-1}\gamma_n\!\left(\left\{u+\frac jM\right\}\right)
                     =A_n(M,\{Mu\}),\qquad Mu\notin\mathbb Z.
\]

\begin{theorem}[Unequal-dilation Stieltjes closure]
\label{mld:dil:full-stieltjes}
Set $A=\log P$, $B=\log Q$, $L_p=\log p$, and $L_q=\log q$.
For all $m,n\ge0$ and the parameters \eqref{mld:dil:parameters},
\begin{align}
 F_{mn}^{p,q}(\theta)={}&
 \sum_{i=0}^m\sum_{j=0}^n\binom mi\binom nj
             (-B)^{m-i}(-A)^{n-j}I_{ij}(\theta)\notag\\
 &+\frac{(-B)^{m+1}-L_p^{m+1}}{m+1}A_n(P,\theta)
  +\frac{(-A)^{n+1}-L_q^{n+1}}{n+1}A_m(Q,1-\theta)\notag\\
 &+\frac{(-B)^{m+1}(-A)^{n+1}}{(m+1)(n+1)}.
 \label{mld:dil:full-formula}
\end{align}
Thus every such product is a finite linear combination of single
Stieltjes values at $\theta$ and $1-\theta$, with coefficients in the
algebra generated over $\mathbb Q$ by the logarithms of $p,q,P,Q$ and
positive integer zeta values.
\end{theorem}

\begin{proof}
Let $C(s,t;\theta)$ be the ordinary translated Hurwitz correlation,
initially for $\Re s,\Re t<1$.  Fourier orthogonality first in
$\Re s,\Re t<0$ gives
\begin{equation}
 \int_0^1\zeta(s,\{px\})\zeta(t,\{qx+a\})\,dx
                   =Q^{s-1}P^{t-1}C(s,t;\theta).
 \label{mld:dil:frequency}
\end{equation}
Indeed the surviving frequency indices satisfy $pu+qv=0$, hence
$u=Q\ell$, $v=-P\ell$; their translation factor is
$e^{-2\pi i\ell Pa}$.  The two sides have the same meromorphic
continuation after local endpoint subtraction.  The common divisor $d$
does not occur in \eqref{mld:dil:frequency}, because ordinary integration is
unchanged by a common covering.

Put $\mathcal G(z,x)=\sum_{n\ge0}\gamma_n(x)z^n/n!$ and
$\mathcal I(z,w;\theta)=\sum_{m,n\ge0}I_{mn}(\theta)z^mw^n/(m!n!)$.
For distinct singular points the distributional Laurent series gives
\begin{equation}
 C(1-z,1-w;\theta)
 =-\frac1{zw}+\frac{\mathcal G(w,\theta)}z
                +\frac{\mathcal G(z,1-\theta)}w
                +\mathcal I(z,w;\theta).
 \label{mld:dil:C-Laurent}
\end{equation}
For example, the coefficient $-1/(zw)$ is the correlation of
$\delta_0-1$ with its nontrivial translate.  Omitting it would change
the final constant in \eqref{mld:dil:full-formula}.

Multiply \eqref{mld:dil:C-Laurent} by $e^{-Bz-Aw}$ and extract
$z^mw^n/(m!n!)$.  This computes the correlation of the \emph{true
pullbacks} $U_pG_m$ and $G_n(qx+a)$.  Its value is the double sum in
\eqref{mld:dil:full-formula}, plus
\[
 \frac{(-B)^{m+1}}{m+1}A_n(P,\theta)
 +\frac{(-A)^{n+1}}{n+1}A_m(Q,1-\theta)
 +\frac{(-B)^{m+1}(-A)^{n+1}}{(m+1)(n+1)}.
\]
The $r=0$ dilation law changes the first factor by
$-L_p^{m+1}\Delta_p/(p(m+1))$.  Pairing this with the other factor
gives $-L_p^{m+1}A_n(P,\theta)/(m+1)$: the grid values form $d$
copies of a complete reduced $P$-grid.  The second factor contributes
$-L_q^{n+1}A_m(Q,1-\theta)/(n+1)$ similarly.  The product of the two
grid distributions vanishes, because their supports are disjoint.
Their sum is \eqref{mld:dil:full-formula}.  The product and all coefficient
extractions can be justified using a partition of unity separating the
two singular grids: on each singular neighborhood the other factor is
smooth, and elsewhere both factors are smooth.
\end{proof}

For example, if $p=2$, $q=1$, $\theta=\{2a\}\in(0,1)$ and
$L=\log2$, \eqref{mld:dil:full-formula} gives the explicit new family member
\begin{align}
 \operatorname{FP}_x\int_0^1\gamma_1(\{2x\})\gamma_0(\{x+a\})\,dx
 ={}&\tfrac12\gamma_2(\theta)+\gamma_2(1-\theta)
        -L\gamma_1(1-\theta)\notag\\
 &+\left(\zeta(2)-\tfrac12L^2\right)\gamma_0(\theta)
        +\zeta(2)\gamma_0(1-\theta)-\zeta(3)-\tfrac12L^3.
 \label{mld:dil:mixed-example}
\end{align}

\subsection{Every mixed parameter derivative}
\begin{theorem}[Dilated differentiation and product integration]
\label{mld:dil:all-derivatives}
Let $N=r+k$ with $r,k,m,n\ge0$, and differentiate
$F_{mn}^{p,q}$ in its displayed argument $\theta$.  Then
\begin{align}
 &\operatorname{FP}_x\int_0^1
       \gamma_m^{(r)}(\{px\})\gamma_n^{(k)}(\{qx+a\})\,dx\notag\\
 &=P^kQ^r\biggl\{
       (-1)^r\frac{d^N}{d\theta^N}F_{mn}^{p,q}(\theta)
       +(-1)^r b_{r,m}(L_p)\partial_\theta^N A_n(P,\theta)\notag\\
 &\hspace{40mm}
       +(-1)^k b_{k,n}(L_q)
          \left.\partial_t^N A_m(Q,t)\right|_{t=1-\theta}
               \biggr\}.
 \label{mld:dil:all-derivative-formula}
\end{align}
The derivatives on the left are derivatives in the \emph{special-function
argument}, not derivatives of the fractional-part map.
\end{theorem}
\begin{proof}
Rewrite \eqref{mld:dil:commutator} as
\[
 X_{p,m,r}=p^{-r}D^rX_{p,m,0}
          +p^{-r-1}b_{r,m}(L_p)\Delta_p^{(r)},
\]
and use the translated counterpart for the second factor.  Periodic
integration by parts in the first product yields
\[
 p^{-r}q^{-k}\int(D^rX_{p,m,0})(D^kX_{q,n,0}(\,\cdot+a/q))
 =(-1)^r(q/p)^r\partial_a^NF_{mn}^{p,q}(\{Pa\}).
\]
Since $\partial_a=P\partial_\theta$ on each noncollision interval,
this is $(-1)^rP^kQ^r\partial_\theta^NF$.
The first grid contact pairs with the smooth second factor and gives
\[
 (-1)^r(q/p)^r b_{r,m}(L_p)\frac1p
       \sum_{j=0}^{p-1}\gamma_n^{(N)}(\{qj/p+a\})
 =(-1)^rP^kQ^r b_{r,m}(L_p)\partial_\theta^NA_n(P,\theta).
\]
The second contact is the third term in
\eqref{mld:dil:all-derivative-formula}.  The grid-grid term vanishes by
disjointness.  The same separated partition of unity used in the
preceding proof makes these integrations by parts valid.
\end{proof}

\subsection{A particularly short polygamma identity}
Put $H_0=0$, $H_j=\sum_{\ell=1}^j1/\ell$ for $j\ge1$, and
\[
 K=\log\operatorname{lcm}(p,q).
\]
\begin{corollary}[All unequal-dilation polygamma correlations]
\label{mld:dil:polygamma}
For $N=r+k\ge1$,
\begin{align}
 &\operatorname{FP}_x\int_0^1
         \psi^{(r)}(\{px\})\psi^{(k)}(\{qx+a\})\,dx\notag\\
 &=N!P^kQ^r\biggl\{
 (-1)^{k+1}\bigl[\zeta'(N+1,\theta)
                 +(H_N-H_r+K)\zeta(N+1,\theta)\bigr]\notag\\
 &\hspace{24mm}+(-1)^{r+1}\bigl[\zeta'(N+1,1-\theta)
                 +(H_N-H_k+K)\zeta(N+1,1-\theta)\bigr]\biggr\}.
 \label{mld:dil:polygamma-formula}
\end{align}
The two primes on the right denote first derivatives in the spectral
variable of Hurwitz zeta.  The $r=k=0$ case is
\begin{align}
 &\operatorname{FP}_x\int_0^1\psi(\{px\})\psi(\{qx+a\})\,dx\notag\\
 &=\gamma_1(\theta)+\gamma_1(1-\theta)
       +K\{\psi(\theta)+\psi(1-\theta)\}
       -\frac{\pi^2}{3}-K^2+\log p\log q.
 \label{mld:dil:digamma-formula}
\end{align}
\end{corollary}
\begin{proof}
Apply \eqref{mld:dil:full-formula} with $m=n=0$.  Its logarithmic
coefficient is $K=L_p+B=L_q+A$, and its constant simplifies using
$AB-K(A+B)=-K^2+L_pL_q$.  This proves
\eqref{mld:dil:digamma-formula}.  In
\eqref{mld:dil:all-derivative-formula}, use $b_{r,0}=H_r$,
$\gamma_0=-\psi$, and the pointwise identities
\[
 \psi^{(N)}(t)=(-1)^{N+1}N!\zeta(N+1,t),\qquad
 \gamma_1^{(N)}(t)=(-1)^{N+1}N!
       \{\zeta'(N+1,t)+H_N\zeta(N+1,t)\}.
\]
The two reflection signs are $(-1)^r$ and $(-1)^k$ before this
conversion.  Simplification gives \eqref{mld:dil:polygamma-formula}.
\end{proof}

The least common multiple in the logarithm is forced by the local
normalization.  For example, a common dilation $p=q=d$ leaves the
ordinary spectral correlation unchanged, but its digamma finite part
changes by $\log d\{\psi(\theta)+\psi(1-\theta)\}$.
Discarding this term would confuse transported finite parts with
ambient-coordinate finite parts.

\begin{corollary}[Odd total derivative order at a half shift]
If $\{Pa\}=1/2$ and $N=r+k$ is odd, then
\begin{align}
 &\operatorname{FP}_x\int_0^1
       \psi^{(r)}(\{px\})\psi^{(k)}(\{qx+a\})\,dx\notag\\
 &=N!P^kQ^r(-1)^{k+1}(H_k-H_r)
                       (2^{N+1}-1)\zeta(N+1).
 \label{mld:dil:half-odd}
\end{align}
In particular every such value is a rational multiple of
$\pi^{N+1}$.  One unequal-dilation example is
\begin{equation}
 \operatorname{FP}_x\int_0^1
       \psi(\{2x\})\psi'(\{3x+1/4\})\,dx=\pi^2.
 \label{mld:dil:pi2-example}
\end{equation}
\end{corollary}
\begin{proof}
When $N$ is odd the two signs in \eqref{mld:dil:polygamma-formula} are
opposite.  At $\theta=1/2$ the spectral derivatives, the $H_N$ terms,
and the $K$ terms cancel pairwise.  The surviving harmonic difference
multiplies $\zeta(N+1,1/2)=(2^{N+1}-1)\zeta(N+1)$.
Since $N+1$ is even, Euler's even-zeta formula gives the last assertion.
\end{proof}

The same unequal pair at zero derivative order gives a different exact
identity, now involving the ordinary first Stieltjes constant:
\begin{align}
 &\operatorname{FP}_x\int_0^1
       \psi(\{2x\})\psi(\{3x+1/4\})\,dx\notag\\
 &=2\gamma_1-2\gamma\log24-\frac{\pi^2}{3}
       -7\log^22-5\log2\log3-\log^23.
 \label{mld:dil:half-digamma-example}
\end{align}
This follows from \eqref{mld:dil:digamma-formula},
$\psi(1/2)=-\gamma-2\log2$ and
$\gamma_1(1/2)=\gamma_1-2\gamma\log2-\log^22$.

\subsection{Ordinary convergent-integral realization}
Every identity in Corollary~\ref{mld:dil:polygamma} can be written without
a finite-part symbol.  Sort the two disjoint singular grids as
$0=x_0<x_1<\cdots<x_{M-1}<1$, put $x_M=1$, and set
$\ell_i=x_{i+1}-x_i$.  On each interval use the ordinary branches
\[
 f_i(t)=\psi^{(r)}(\alpha_i+pt)
        \psi^{(k)}(\beta_i+qt),\qquad0<t<\ell_i,
\]
where $\alpha_i=\{px_i\}$ and $\beta_i=\{qx_i+a\}$, with the
right-hand branch chosen at zero.  Exactly one of $\alpha_i,\beta_i$
is zero.  If $\alpha_i=0$, set $v_i=r$ and
\[
 a_{ij}=\frac{(-1)^{r+1}r!q^j}{p^{r+1}j!}
                       \psi^{(k+j)}(\beta_i),\quad0\le j\le r.
\]
If $\beta_i=0$, set $v_i=k$ and
\[
 a_{ij}=\frac{(-1)^{k+1}k!p^j}{q^{k+1}j!}
                       \psi^{(r+j)}(\alpha_i),\quad0\le j\le k.
\]
Then the left side of \eqref{mld:dil:polygamma-formula}, including its
$N=0$ counterpart, equals the ordinary absolutely convergent expression
\begin{equation}
 \sum_{i=0}^{M-1}\left[
  \int_0^{\ell_i}\left\{f_i(t)
               -\sum_{j=0}^{v_i}a_{ij}t^{j-v_i-1}\right\}\,dt
  +a_{i,v_i}\log\ell_i
  +\sum_{j=0}^{v_i-1}\frac{a_{ij}\ell_i^{j-v_i}}{j-v_i}
                           \right].
 \label{mld:dil:ordinary-realization}
\end{equation}
This follows by integrating the displayed singular monomials explicitly.
The subtracted integrands are bounded at their left endpoints.  At the
right endpoint the branch approaching one is used, so no singularity
is introduced there.

The independent numerical implementation uses precisely
\eqref{mld:dil:ordinary-realization} on the left and
\eqref{mld:dil:polygamma-formula}--\eqref{mld:dil:digamma-formula} on the right.
Nine cases include unequal dilations, a nontrivial common divisor, and
derivative orders up to $r+k=3$.  At $42$ decimal working digits the
largest relative discrepancy is below $3.5\times10^{-39}$.  These are
floating-point consistency checks of formulas proved above, not certified
error enclosures or a substitute for the analytic argument.
Three further independent subtracted-integral checks use positive
Stieltjes indices $(m,n)=(1,0),(0,1),(1,1)$, the last with
$(p,q)=(2,3)$.  At $28$ working digits their largest relative discrepancy
is below $9\times10^{-29}$.
